% Portable theorem guard: define only what the host preamble is missing, so this
% file drops into any of our manuscripts and re-running it is a no-op.
\makeatletter
\@ifundefined{theorem}{\newtheorem{theorem}{Theorem}}{}
\@ifundefined{lemma}{\newtheorem{lemma}[theorem]{Lemma}}{}
\@ifundefined{proposition}{\newtheorem{proposition}[theorem]{Proposition}}{}
\@ifundefined{corollary}{\newtheorem{corollary}[theorem]{Corollary}}{}
\@ifundefined{definition}{\newtheorem{definition}[theorem]{Definition}}{}
\@ifundefined{remark}{\newtheorem{remark}[theorem]{Remark}}{}
\makeatother
% BT1154 projective 15-sector refinement insert.

\begin{theorem}[Projective $15$-sector refinement]
\label{thm:projective-fifteen-refinement}
The $15$-sector suggested by Remark~\ref{rem:projective-fifteen-sector} is real,
but not as the naive support-incidence module.  The $W(3,3)$ collinearity matrix
has negative projector
\[
  P_-={(A-12I)(A-2I)\over 96},\qquad \operatorname{rank}P_-=15.
\]
Mapping a projective $\mathbb F_3$ point to its nonzero coordinate-support gives
15 support fibers indexed by $\mathrm{PG}(3,2)$, but their projection into
$P_-$ has rank $5$, not $15$.  Thus the naive support-fiber bridge is refuted.
The surviving bridge is the Boolean/Clifford character completion
\[
  16=1+15,
\]
where the extra line is the scalar/vacuum empty mask and the $15$ nonempty masks
are the projective sector.  The grade-four pseudoscalar is inside the $15$,
not the extra line.
\end{theorem}
